\documentclass[11pt]{article}
\usepackage{amsmath,amssymb,amsthm}
\usepackage[margin=1in]{geometry}
\newtheorem{theorem}{Theorem}
\newtheorem{lemma}[theorem]{Lemma}
\newtheorem{corollary}[theorem]{Corollary}
\theoremstyle{definition}
\newtheorem{definition}[theorem]{Definition}
\title{Seymour's second neighbourhood conjecture at order $n=2\delta+3$:\\ a reduction to complete minimum-outdegree vertices}
\author{Working note, winter 2026 (status: PROVED items below, case still open)}
\date{17 November 2026}
\begin{document}
\maketitle

\section*{Setting}
$D$ is a finite oriented graph on $n$ vertices with minimum outdegree $\delta$. For a vertex $v$ write $N_1(v)$ for its
out-neighbourhood, $d(v)=|N_1(v)|$, $j_v=d(v)-\delta$, $N_2(v)$ for the set of vertices at out-distance exactly $2$,
$R(v)=\{v\}\cup N_1(v)\cup N_2(v)$ and $U(v)=V\setminus R(v)$. For a target $x$ let $T(x)$ be the set of vertices
$y\neq x$ with $y\not\to x$, $q_x=|T(x)|-\delta$, $P(x)=\{v: x\in U(v)\}$ and $p_x=|P(x)|$. Let $M$ be the set of
non-adjacent pairs and $\sigma_v=d(v)-1-|N_2(v)|$. $D$ has no Seymour vertex iff $\sigma_v\ge 0$ for every $v$.
Brukhman (arXiv:2608.11530) proved the conjecture when $n=2\delta+2$; Fisher's theorem covers $n=2\delta+1$.
Throughout, $n=2\delta+3$, the first open dense order.

\begin{lemma}[exact capacity]\label{cap}
If $p=p_x>0$ then $N_1(v)\subseteq T(x)$ for every $v\in P(x)$, and
\[ \tfrac12\,p\,(2q_x-1-p) \;=\; J_P+m_P+r_x, \]
where $J_P=\sum_{v\in P(x)}j_v$, $m_P$ is the number of non-adjacent pairs inside $P(x)$, and $r_x$ is the number of
pairs $(v,t)$ with $v\in P(x)$, $t\in T(x)\setminus P(x)$, $v\not\to t$.
\end{lemma}
\begin{proof}
If $v\in P(x)$ and $v\to y$ then $y\ne x$ and $y\not\to x$, so $y\in T(x)$; also $v\in T(x)$. Every arc leaving
$P(x)$ therefore ends in $T(x)$. Counting these arcs exactly,
$\sum_{v\in P}(\delta+j_v)=\binom p2-m_P+p(|T(x)|-p)-r_x$ with $|T(x)|=\delta+q_x$. Rearranging gives the identity.
(This is Brukhman's Lemma 2.1 with the slack terms kept.)
\end{proof}

\begin{lemma}[identity]\label{id}
If $n=2\delta+3$ then $\sum_x (p_x-2q_x+1)=\sum_v \sigma_v$, and $q_x=\deg_M(x)+j_x$.
\end{lemma}
\begin{proof}
$\sum_x p_x=\sum_v|U(v)|=\sum_v(n-1-d(v)-|N_2(v)|)$ and $\sum_x q_x=n(n-1-\delta)-\sum_v d(v)$, so the left side equals
$n(2\delta+2-n)+\sum_v(d(v)-|N_2(v)|)=\sum_v\sigma_v$. For the second claim, $|T(x)|=n-1-d^-(x)=\deg_M(x)+d(x)$.
\end{proof}

Write $D_x=2q_x-1-p_x$. Call $x$ \emph{tight} if $p_x>0$ and $D_x=0$, and let
$Z=\{x: q_x=0\}$, the vertices with no missing edge and outdegree $\delta$.

\begin{lemma}[tight targets are covered]\label{tight}
If $x$ is tight and $v\in P(x)$, then $T(x)\subseteq R(v)$.
\end{lemma}
\begin{proof}
By Lemma~\ref{cap}, $J_P=m_P=r_x=0$: $P(x)$ is a tournament on $p=2q-1$ vertices ($q=q_x$), all of outdegree $\delta$,
and each $v\in P(x)$ dominates $T(x)\setminus P(x)$, which has $\delta-q+1$ elements. So $v$ has exactly $q-1$
out-neighbours inside $P(x)$, and $P(x)$ is a regular tournament. If $w\in P(x)$ and $w\to v$, then $w$ cannot beat
all $q-1$ out-neighbours of $v$ in $P(x)$, since $w$ would then have outdegree at least $q$ inside $P(x)$; hence some
out-neighbour of $v$ beats $w$ and $w\in N_2(v)$. Thus $P(x)\cup(T(x)\setminus P(x))\subseteq R(v)$.
\end{proof}

\begin{lemma}[at most one covered target]\label{one}
For every vertex $v$ of any oriented graph, at most one $x\in U(v)$ satisfies $T(x)\subseteq R(v)$.
\end{lemma}
\begin{proof}
If $T(x)\subseteq R(v)$ and $y\in U(v)\setminus\{x\}$, then $y\notin T(x)$ and $y\neq x$, so $y\to x$. Two such
targets $x_1,x_2$ would give $x_1\to x_2\to x_1$.
\end{proof}

\begin{theorem}\label{main}
Let $D$ be a counterexample with $n=2\delta+3$. Let $N$ be the set of targets with $p_x>0$ and $D_x\ge 1$, and $E$ the
set of $x\notin Z$ with $p_x=0$. Then
\[ |Z| \;=\; \sum_v\sigma_v+\sum_{x\in E}(2q_x-1)+\sum_{x\in N}D_x \;\ge\; |N| \;\ge\; \max_v |U(v)|-1 \;\ge\; 2 .\]
In particular $D$ has at least two vertices that have no missing edge and outdegree $\delta$, and every vertex has at
most one tight target.
\end{theorem}
\begin{proof}
For $x\notin Z$ with $p_x>0$, Lemma~\ref{cap} gives $D_x\ge0$; for $x\in Z$, $q_x=0$ forces $p_x=0$ by
Lemma~\ref{cap} and $D_x=-1$. The equation is Lemma~\ref{id} rewritten. Every $x\in U(v)$ has $p_x>0$, and by
Lemmas~\ref{tight} and~\ref{one} at most one of them is tight, so at least $|U(v)|-1$ of them lie in $N$. A vertex $v$
of outdegree $\delta$ has $|U(v)|=n-1-\delta-|N_2(v)|\ge 3$.
\end{proof}

\begin{corollary}
If a counterexample with $n=2\delta+3$ has no vertex that is complete (adjacent to all others) of outdegree $\delta$,
or has only one, it does not exist. In particular SSNC holds for every oriented graph with $n=2\delta+3$ in which at
most one vertex of outdegree $\delta$ is adjacent to all other vertices.
\end{corollary}

\begin{lemma}[targets of a complete vertex]\label{zt}
Let $z\in Z$ and $k=|U(z)|$. Then $N^-(z)=N_2(z)\cup U(z)$; for every $y\in U(z)$ all in-neighbours of $y$ lie in
$N^-(z)$, $T(y)=\{z\}\cup N_1(z)\cup X_y$ with $X_y\subseteq N^-(z)\setminus\{y\}$ and $|X_y|=q_y-1$; $y$ is covered
for $z$ iff $X_y\cap U(z)=\emptyset$; at most one $y\in U(z)$ has $q_y=1$; and
$\sum_{y\in U(z)}q_y\ge k(k+1)/2$.
\end{lemma}
\begin{proof}
$z$ is adjacent to every vertex and has $\delta$ out-neighbours, so its $\delta+2$ in-neighbours are exactly the vertices
outside $R(z)$ or in $N_2(z)$. For $y\in U(z)$ no vertex of $\{z\}\cup N_1(z)$ sends an arc to $y$, which gives the
description of $T(y)$; $q_y=1$ means $X_y=\emptyset$, so $y$ is covered, and Lemma~\ref{one} allows one such $y$. The
in-arcs of the vertices of $U(z)$ come from the set $N^-(z)$ of size $\delta+2$: at most $\binom k2$ from inside $U(z)$
and at most $k(\delta+2-k)$ from $N_2(z)$. Hence $\sum_{y\in U(z)}(\delta+2-q_y)\le\binom k2+k(\delta+2-k)$, which
rearranges to the bound.
\end{proof}

\section*{Status}
The remaining case, $|Z|\ge 2$, is open. Each $z\in Z$ satisfies $N^-(z)=N_2(z)\cup U(z)$, $|N^-(z)|=\delta+2$, and has
at least two non-tight targets in $U(z)\subseteq N^-(z)$. Closing the case needs a charging of $Z$ to the deficiencies
$D_x$, $x\in N$.
\end{document}
